\documentclass[../DoAn.tex]{subfiles}
\begin{document}


\section{Đặt vấn đề}
\label{sec:dvd}

Trong môi trường giáo dục đại học hiện nay, đặc biệt tại các trường có quy mô lớn và đào tạo chuyên sâu như \textbf{Trường Đại học Bách Khoa Hà Nội}, việc tra cứu thông tin về quy chế, quy định đào tạo là nhu cầu thiết yếu của sinh viên. Các câu hỏi thường gặp bao gồm:

\begin{itemize}
    \item Điều kiện tốt nghiệp, yêu cầu ngoại ngữ;
    \item Quy định về học bổng, miễn giảm học phí;
    \item Thủ tục xin nghỉ học, bảo lưu, đăng ký môn học;
    \item Quy chế thi cử, đánh giá kết quả học tập;
    \item Hướng dẫn chuyển ngành, chuyển tiếp.
\end{itemize}

Với hệ thống văn bản quy định phức tạp gồm hàng chục quyết định, thông tư và hướng dẫn, sinh viên thường gặp khó khăn trong việc tra cứu thông tin chính xác. Việc đọc và tìm kiếm thủ công trong các văn bản pháp lý dài không chỉ tốn thời gian mà còn dễ bỏ sót thông tin quan trọng.


\section{Các giải pháp hiện tại và hạn chế}
\label{sec:giaiphap}

Hiện nay, việc tra cứu thông tin quy chế đào tạo của sinh viên chủ yếu dựa vào các phương thức sau:

\begin{itemize}
    \item \textbf{Tra cứu thủ công:} Sinh viên tự tìm kiếm trong các văn bản PDF được đăng tải trên website trường, đòi hỏi thời gian và kỹ năng tìm kiếm;
    
    \item \textbf{Hỏi cán bộ phòng đào tạo:} Gửi email hoặc đến trực tiếp để được giải đáp, nhưng thường phải chờ đợi và số lượng nhân viên có hạn;
    
    \item \textbf{Hỏi đáp trên diễn đàn:} Thông tin có thể không chính xác hoặc lỗi thời do sinh viên trả lời thay vì nguồn chính thức;
    
    \item \textbf{Chatbot đơn giản:} Một số chatbot dựa trên từ khóa (rule-based) không có khả năng hiểu ngữ cảnh và trả lời linh hoạt.
\end{itemize}

Các hạn chế chính của giải pháp hiện tại:
\begin{itemize}
    \item Chưa tận dụng công nghệ AI và NLP hiện đại;
    \item Không có khả năng trích dẫn nguồn cụ thể (Điều, Khoản);
    \item Thông tin trả lời thiếu độ tin cậy và khó xác minh.
\end{itemize}


\section{Mục tiêu và định hướng giải pháp}
Trước nhu cầu cấp thiết về việc hỗ trợ sinh viên tra cứu quy chế đào tạo một cách hiệu quả, đồ án đề xuất xây dựng hệ thống RAG (Retrieval-Augmented Generation) Chatbot với các mục tiêu sau:

\begin{itemize}
    \item \textbf{Hỏi-đáp bằng ngôn ngữ tự nhiên:} Sinh viên có thể đặt câu hỏi bằng tiếng Việt tự nhiên, không cần biết cấu trúc văn bản;
    
    \item \textbf{Trích dẫn nguồn chính xác:} Mỗi câu trả lời kèm theo Điều, Khoản cụ thể từ văn bản quy chế;

    \item \textbf{Xử lý đa nguồn tài liệu:} Tổng hợp thông tin từ 16+ văn bản quy định đào tạo;

    \item \textbf{Độ chính xác cao:} Đạt accuracy 99.5\%+ trong xử lý tiếng Việt.
\end{itemize}

\section{Đóng góp của đồ án}
\label{sec:donggop}

Đồ án có năm đóng góp chính như sau:

\begin{enumerate}
    \item \textbf{Xây dựng pipeline xử lý PDF đa phương thức:} Kết hợp Docling, PyMuPDF4LLM và olmOCR để xử lý các văn bản PDF với font encoding phức tạp, đảm bảo độ chính xác 99.5\%+ cho tiếng Việt;

    \item \textbf{Phát triển 4 chiến lược chunking chuyên biệt:} Thiết kế OlmOcrLegalChunker, RecursiveChunker, STSVChunker và ArticleLevelLegalChunker với khả năng nhận diện cấu trúc Chương-Điều-Khoản-Điểm, bảo toàn ngữ cảnh và metadata;

    \item \textbf{Triển khai Hybrid Search với Dual Embedding:} Kết hợp BGE-M3 và E5-Multilingual lưu trên Qdrant Vector Database, kết hợp với Elasticsearch BM25 qua Reciprocal Rank Fusion đạt Hit@5 = 86.4\%;

    \item \textbf{Xây dựng hệ thống RAG 8 tầng hoàn chỉnh:} Tích hợp Query Router, Query Reflection, Cross-encoder Reranking, Gemini 2.5 Flash, Self Evaluation, Tavily fallback và MongoDB Memory;

    \item \textbf{Đánh giá quy mô lớn:} Xây dựng framework đánh giá với 25 bộ dữ liệu, 2.390 câu hỏi, so sánh 6 phương pháp retrieval.
\end{enumerate}


\section{Bố cục đồ án}

Phần còn lại của báo cáo đồ án tốt nghiệp này được tổ chức như sau.

Chương 2 trình bày cơ sở lý thuyết về hệ thống RAG, bao gồm kiến thức nền tảng về NLP, embedding (BGE-M3, E5-Multilingual), vector store (Qdrant), và các kỹ thuật chunking cho văn bản pháp lý. Chương này cũng khảo sát các nghiên cứu liên quan và phân tích các công nghệ sử dụng.

Chương 3 trình bày kiến trúc tổng thể của hệ thống RAG Chatbot với 8 tầng, mô tả chi tiết pipeline thu thập dữ liệu đa nguồn, các chiến lược chunking chuyên biệt, và luồng xử lý từ query đến response.

Chương 4 đi sâu vào việc cài đặt và triển khai toàn bộ các module: xử lý dữ liệu và chunking, Embedding Layer (BGE-M3, E5), Hybrid Retrieval (Qdrant + Elasticsearch + RRF), Reranking (BGE cross-encoder), Query Router \& Reflection, Chat Model (GeminiLLM) và Self Evaluation, Tavily fallback, MongoDB Memory, FastAPI Backend và Pipeline Integration.

Chương 5 tập trung vào việc đánh giá hệ thống trên 2.390 câu hỏi từ 25 bộ dữ liệu, so sánh 6 phương pháp retrieval với các metrics Hit@K, Precision@K, Recall@K, MRR.

Cuối cùng, chương 6 đưa ra kết luận và hướng phát triển của đồ án, tổng kết các đóng góp chính, nhìn nhận hạn chế và đề xuất các cải tiến bao gồm query caching, metadata-aware retrieval, và nâng cấp Agentic RAG.

\end{document}
